\documentclass[11pt]{article}
\usepackage[margin=1in]{geometry}
\usepackage[T1]{fontenc}
\usepackage{lmodern,microtype,hyperref,enumitem,amsmath,amssymb,xcolor,booktabs}
\usepackage[numbers,sort&compress]{natbib}
\hypersetup{colorlinks=true,urlcolor=blue,citecolor=blue,linkcolor=blue}
\title{Q14: How Are ``Concepts'' Represented---Do They Live in Features?}
\author{DLS 2026 --- module author: Kukla}
\date{October 2026}
\usepackage{url}
\begin{document}
\maketitle

\section*{Learning objectives}
\begin{itemize}[leftmargin=1.4em]
\item Explain how neural networks represent ``concepts'' in activation space.
\item Define superposition and why it complicates one-neuron-one-concept intuitions.
\item Distinguish a model's weights (long-term knowledge) from its context (working memory).
\end{itemize}

\begin{quote}\small\textbf{Companion video:}\begin{itemize}\setlength\itemsep{0pt}\item \emph{Reading AI's Mind - Mechanistic Interpretability Explained} --- bycloud (9:21). \url{https://www.youtube.com/watch?v=Mhp8vpOksWw}\end{itemize}\end{quote}

\section{The question}
Student response \#27 asks how ``concepts'' are maintained and generated, whether they live in features,
and how models reason and keep memory. The modern answer from interpretability research is specific:
concepts are \textbf{directions in activation space}, often many more of them than there are neurons, and
``memory'' is really two different things that students should not conflate---weights and context.

\section{Concepts as directions, not neurons}
Early intuition hoped each neuron encoded one human-meaningful concept. Reality is messier and more
elegant: a concept typically corresponds to a \emph{direction} (a linear combination of neuron
activations), and a single neuron participates in many concepts \citep{olah2020zoom}. This
\textbf{linear representation hypothesis}---that features are directions and that meaningful operations are
approximately linear in that space---is supported by findings like vector arithmetic on word embeddings and
by probes that linearly decode attributes from hidden states. Reasoning, on this view, is the model moving
and combining these feature directions across layers via attention and the MLP.

\section{Superposition: more features than neurons}
Why not just read off neurons? Because networks pack \emph{more features than they have dimensions}, a
phenomenon called \textbf{superposition} \citep{elhage2022superposition}. When features are sparse (only a
few active at once), a layer can represent many of them as nearly-orthogonal directions in a lower-
dimensional space, tolerating small interference. Superposition explains \textbf{polysemantic neurons}
(single units that fire for unrelated inputs) and is the central obstacle to naive interpretability. The
response---\textbf{sparse autoencoders}---learns an overcomplete, sparse dictionary that disentangles the
superposed features back into individually meaningful ones \citep{bricken2023monosemanticity}, a major
recent step in reading concepts out of a model (developed further in the interpretability module Q27).

\section{Where knowledge is stored: the MLP layers}
Factual, associative knowledge lives disproportionately in the \textbf{MLP (feed-forward) sublayers},
which behave like key--value memories: a key pattern in the residual stream retrieves an associated value
written back into it \citep{geva2021kv}. Causal-tracing work can localize \emph{where} a specific fact is
stored and even edit it \citep{meng2022rome}. Attention, by contrast, is mainly about \emph{moving}
information between positions; the MLPs are where much of the ``what the model knows'' resides.

\section{A worked intuition: the king--queen direction}
The classic demonstration that concepts are directions is word-vector arithmetic: in a well-trained
embedding space, the vector for ``king'' minus ``man'' plus ``woman'' lands near ``queen.'' The feature
``royalty'' and the feature ``gender'' behave like (approximately) independent directions you can add and
subtract. Modern mechanistic work finds the same structure deeper in the network: a ``this text is in
French'' direction, a ``sentiment is positive'' direction, a ``the subject is plural'' direction---each
recoverable by a linear probe, each manipulable by adding a scaled vector to the residual stream to
\emph{steer} the model's behavior. This is why the linear-representation picture is more than a metaphor:
it makes testable predictions (probes should decode the attribute; adding the direction should change the
output) that largely hold. It also clarifies ``how models reason'': reasoning is, in large part, the
network reading some feature directions and writing others across layers, routed by attention.

\section{When the clean picture bends}
Honesty requires noting the limits. Not every feature is cleanly linear or cleanly separable; superposition
means directions interfere, and some features appear to be represented in more complex, possibly nonlinear
ways that current tools capture poorly. Sparse autoencoders recover many interpretable features but not all,
and which dictionary you learn depends on hyperparameters. So ``concepts are directions'' is the best
current working model---powerful, predictive, and actively used for steering and interpretability---rather
than a complete theory. Students should hold it as a strong, useful approximation, not dogma.

\section{Two kinds of memory}
Finally, the word ``memory'' hides a crucial split. \textbf{Weights} are long-term memory: learned during
training, fixed at inference, holding the concepts and facts above. \textbf{Context} (the tokens in the
window) is short-term working memory: available now, gone when the conversation ends, not written back into
weights. A model does not ``remember'' yesterday's chat unless that text is re-supplied in context or
retrieved (module Q15). Confusing the two is the source of many misconceptions---why context is not memory
is the subject of the next module.

\section{Summary}
Concepts are directions in activation space, not single neurons; superposition lets a network hold more
features than dimensions, and sparse autoencoders can disentangle them. Factual knowledge is concentrated
in MLP key--value memories and can be localized and edited. ``Memory'' splits into weights (long-term,
trained) and context (short-term, working)---a distinction module Q15 builds on.

\bibliographystyle{plainnat}
\bibliography{bibliography}
\end{document}
